
\documentclass{article}
\usepackage{colortbl}
\usepackage{makecell}
\usepackage{multirow}
\usepackage{supertabular}

\begin{document}

\newcounter{utterance}

\centering \large Interaction Transcript for game `dond', experiment `coop\_en', episode 19 with Qwen3{-}32B{-}t0.0.
\vspace{24pt}

{ \footnotesize  \setcounter{utterance}{1}
\setlength{\tabcolsep}{0pt}
\begin{supertabular}{c@{$\;$}|p{.15\linewidth}@{}p{.15\linewidth}p{.15\linewidth}p{.15\linewidth}p{.15\linewidth}p{.15\linewidth}}
    \# & $\;$A & \multicolumn{4}{c}{Game Master} & $\;\:$B\\
    \hline

    \theutterance \stepcounter{utterance}  
    & & \multicolumn{4}{p{0.6\linewidth}}{
        \cellcolor[rgb]{0.9,0.9,0.9}{
            \makecell[{{p{\linewidth}}}]{
                \texttt{\tiny{[P1$\langle$GM]}}
                \texttt{You are playing a cooperative negotiation game in which you have to agree on how to divide a set of items among you and another player.} \\
\\ 
\texttt{Rules:} \\
\texttt{(a) You and the other player are given a set of items. Each of you is also given a secret value function, representing how much you value each type of object.} \\
\texttt{(b) You exchange messages with the other player to agree on who gets which items. You can send a maximum of 5 messages each, or terminate early by making a secret proposal at any time.} \\
\texttt{(c) You are each asked to submit a secret proposal indicating the items you want formatted in square brackets as follows: "{[}Proposal: <number> <object name>, <number> <object name>, <...>{]}"} \\
\texttt{(d) If your proposals are complementary, i.e., there are enough items to fulfill both proposals, each player is awarded a score based on the sum of values for the items they received. Otherwise, both of you get zero points.} \\
\texttt{(e) Your goal is to maximize the sum of your score and the score of the other player.} \\
\\ 
\texttt{Let us start.} \\
\\ 
\texttt{The set of available items is:} \\
\\ 
\texttt{2 hooks, 1 flower, 1 card, 1 eraser, 1 sponge} \\
\\ 
\texttt{Your secret value function is:} \\
\\ 
\texttt{hook: 3, flower: 3, card: 0, eraser: 1, sponge: 0} \\
\\ 
\texttt{IMPORTANT: Your messages, unless it is the secret proposal, are directly transmitted to the other player, so do not include any response to the rules or text announcing your message. To make a secret proposal, use the indicated format. Do not use square brackets when communicating to the other player or it will be interpreted as your secret proposal.} \\
            }
        }
    }
    & \\ \\

    \theutterance \stepcounter{utterance}  
    & \multicolumn{4}{p{0.6\linewidth}}{
        \cellcolor[rgb]{0.9,0.9,0.9}{
            \makecell[{{p{\linewidth}}}]{
                \texttt{\tiny{[P1$\rangle$GM]}}
                \texttt{<think>} \\
\texttt{Okay, let's see. I need to divide these items between me and another player. The items are 2 hooks, 1 flower, 1 card, 1 eraser, 1 sponge. My values are hook:3, flower:3, card:0, eraser:1, sponge:0. The other player has their own secret values, which I don't know. The goal is to maximize the total score for both of us, so we need to find a way to split the items where our proposals don't conflict.} \\
\\ 
\texttt{First, I should figure out what I value most. Hooks and flowers are worth 3 each, eraser is 1. The other items are 0. So I probably want to get as many hooks and flowers as possible. But there are only two hooks and one flower. If I take both hooks and the flower, that's 3+3+3=9 points for me. But then the other player would have to take the card, eraser, and sponge. But maybe the other player values some of those items more. Wait, but the problem is that if our proposals are complementary, meaning we can both get what we want without overlapping. So I need to propose a division that allows the other player to also get their desired items.} \\
\\ 
\texttt{But since I don't know their values, I have to make some assumptions. Maybe the other player also values some of the items I don't. For example, maybe they value the card or sponge. But since I can't know, I need to communicate. The rules say we can exchange up to 5 messages each. So I should start by sending a message to the other player to discuss possible allocations.} \\
\\ 
\texttt{Let me think. If I propose taking both hooks and the flower, that's three items. The other player would need to take the remaining items: card, eraser, sponge. But maybe they don't want those. Alternatively, maybe I can offer to give up one hook in exchange for something else they value. For example, if they value the eraser more, maybe I can take one hook and the flower, and they take the other hook and the eraser. But I don't know their values. Alternatively, maybe they want the flower. If I take one hook and the flower, and they take the other hook and maybe the eraser. But again, without knowing their values, it's tricky.} \\
\\ 
\texttt{Wait, the key is} \\
            }
        }
    }
    & & \\ \\

    \theutterance \stepcounter{utterance}  
    & & & \multicolumn{2}{p{0.3\linewidth}}{
        \cellcolor[rgb]{0.9,0.9,0.9}{
            \makecell[{{p{\linewidth}}}]{
                \texttt{\tiny{[GM$|$GM]}}
                \texttt{continue} \\
            }
        }
    }
    & & \\ \\

    \theutterance \stepcounter{utterance}  
    & & & \multicolumn{2}{p{0.3\linewidth}}{
        \cellcolor[rgb]{0.9,0.9,0.9}{
            \makecell[{{p{\linewidth}}}]{
                \texttt{\tiny{[GM$|$GM]}}
                \texttt{<think>} \\
\texttt{Okay, let's see. I need to divide these items between me and another player. The items are 2 hooks, 1 flower, 1 card, 1 eraser, 1 sponge. My values are hook:3, flower:3, card:0, eraser:1, sponge:0. The other player has their own secret values, which I don't know. The goal is to maximize the total score for both of us, so we need to find a way to split the items where our proposals don't conflict.} \\
\\ 
\texttt{First, I should figure out what I value most. Hooks and flowers are worth 3 each, eraser is 1. The other items are 0. So I probably want to get as many hooks and flowers as possible. But there are only two hooks and one flower. If I take both hooks and the flower, that's 3+3+3=9 points for me. But then the other player would have to take the card, eraser, and sponge. But maybe the other player values some of those items more. Wait, but the problem is that if our proposals are complementary, meaning we can both get what we want without overlapping. So I need to propose a division that allows the other player to also get their desired items.} \\
\\ 
\texttt{But since I don't know their values, I have to make some assumptions. Maybe the other player also values some of the items I don't. For example, maybe they value the card or sponge. But since I can't know, I need to communicate. The rules say we can exchange up to 5 messages each. So I should start by sending a message to the other player to discuss possible allocations.} \\
\\ 
\texttt{Let me think. If I propose taking both hooks and the flower, that's three items. The other player would need to take the remaining items: card, eraser, sponge. But maybe they don't want those. Alternatively, maybe I can offer to give up one hook in exchange for something else they value. For example, if they value the eraser more, maybe I can take one hook and the flower, and they take the other hook and the eraser. But I don't know their values. Alternatively, maybe they want the flower. If I take one hook and the flower, and they take the other hook and maybe the eraser. But again, without knowing their values, it's tricky.} \\
\\ 
\texttt{Wait, the key is} \\
            }
        }
    }
    & & \\ \\

    \theutterance \stepcounter{utterance}  
    & & \multicolumn{4}{p{0.6\linewidth}}{
        \cellcolor[rgb]{0.9,0.9,0.9}{
            \makecell[{{p{\linewidth}}}]{
                \texttt{\tiny{[GM$\rangle$P2]}}
                \texttt{You are playing a cooperative negotiation game in which you have to agree on how to divide a set of items among you and another player.} \\
\\ 
\texttt{Rules:} \\
\texttt{(a) You and the other player are given a set of items. Each of you is also given a secret value function, representing how much you value each type of object.} \\
\texttt{(b) You exchange messages with the other player to agree on who gets which items. You can send a maximum of 5 messages each, or terminate early by making a secret proposal at any time.} \\
\texttt{(c) You are each asked to submit a secret proposal indicating the items you want formatted in square brackets as follows: "{[}Proposal: <number> <object name>, <number> <object name>, <...>{]}"} \\
\texttt{(d) If your proposals are complementary, i.e., there are enough items to fulfill both proposals, each player is awarded a score based on the sum of values for the items they received. Otherwise, both of you get zero points.} \\
\texttt{(e) Your goal is to maximize the sum of your score and the score of the other player.} \\
\\ 
\texttt{Let us start.} \\
\\ 
\texttt{The set of available items is:} \\
\\ 
\texttt{2 hooks, 1 flower, 1 card, 1 eraser, 1 sponge} \\
\\ 
\texttt{Your secret value function is:} \\
\\ 
\texttt{hook: 0, flower: 2, card: 1, eraser: 4, sponge: 3} \\
\\ 
\texttt{IMPORTANT: Your messages, unless it is the secret proposal, are directly transmitted to the other player, so do not include any response to the rules or text announcing your message. To make a secret proposal, use the indicated format. Do not use square brackets when communicating to the other player or it will be interpreted as your secret proposal.} \\
\\ 
\\ 
\texttt{<think>} \\
\texttt{Okay, let's see. I need to divide these items between me and another player. The items are 2 hooks, 1 flower, 1 card, 1 eraser, 1 sponge. My values are hook:3, flower:3, card:0, eraser:1, sponge:0. The other player has their own secret values, which I don't know. The goal is to maximize the total score for both of us, so we need to find a way to split the items where our proposals don't conflict.} \\
\\ 
\texttt{First, I should figure out what I value most. Hooks and flowers are worth 3 each, eraser is 1. The other items are 0. So I probably want to get as many hooks and flowers as possible. But there are only two hooks and one flower. If I take both hooks and the flower, that's 3+3+3=9 points for me. But then the other player would have to take the card, eraser, and sponge. But maybe the other player values some of those items more. Wait, but the problem is that if our proposals are complementary, meaning we can both get what we want without overlapping. So I need to propose a division that allows the other player to also get their desired items.} \\
\\ 
\texttt{But since I don't know their values, I have to make some assumptions. Maybe the other player also values some of the items I don't. For example, maybe they value the card or sponge. But since I can't know, I need to communicate. The rules say we can exchange up to 5 messages each. So I should start by sending a message to the other player to discuss possible allocations.} \\
\\ 
\texttt{Let me think. If I propose taking both hooks and the flower, that's three items. The other player would need to take the remaining items: card, eraser, sponge. But maybe they don't want those. Alternatively, maybe I can offer to give up one hook in exchange for something else they value. For example, if they value the eraser more, maybe I can take one hook and the flower, and they take the other hook and the eraser. But I don't know their values. Alternatively, maybe they want the flower. If I take one hook and the flower, and they take the other hook and maybe the eraser. But again, without knowing their values, it's tricky.} \\
\\ 
\texttt{Wait, the key is} \\
            }
        }
    }
    & \\ \\

    \theutterance \stepcounter{utterance}  
    & & & \multicolumn{4}{p{0.6\linewidth}}{
        \cellcolor[rgb]{0.9,0.9,0.9}{
            \makecell[{{p{\linewidth}}}]{
                \texttt{\tiny{[GM$\langle$P2]}}
                \texttt{<think>} \\
\texttt{Okay, let's start by understanding the problem. I need to divide the items between me and another player. The items are 2 hooks, 1 flower, 1 card, 1 eraser, 1 sponge. My value function is hook:3, flower:3, card:0, eraser:1, sponge:0. The other player's values are secret, so I don't know what they want. The goal is to maximize the total score for both of us, which means we need to find a division where our proposals don't conflict. If they do conflict, we both get zero, so it's crucial that our proposals are complementary.} \\
\\ 
\texttt{First, I should figure out what I value most. Hooks and flowers are the highest for me, each worth 3. The eraser is 1, and the others are 0. So ideally, I want to get as many hooks and flowers as possible. But there are only two hooks and one flower. If I take both hooks and the flower, that's 3+3+3=9 points for me. But then the other player would have to take the card, eraser, and sponge. However, maybe the other player values some of those items. For example, if they value the eraser or sponge, they might get some points from those. But if their values are different, maybe they want some of the items I have. But since I don't know their values, I need to communicate.} \\
\\ 
\texttt{The rules say we can exchange up to 5 messages each. So I should start by sending a message to the other player to discuss possible allocations. Let me think about what to propose. Since I don't know their preferences, maybe I can suggest a division that allows them to take items I don't value. For example, if I take the two hooks and the flower, I can offer the other player the eraser and sponge. But the card is also there. Alternatively, maybe I can take one hook and the flower, and let them take the other hook and the eraser. But I need to make sure that the other player can get their desired items without overlapping with mine.} \\
\\ 
\texttt{Wait, the key is to make sure that the proposals are complementary. So if I propose taking certain items, the other player's proposal must not require any of those items. For example, if I take both hooks and the flower, the other player can't take any of} \\
            }
        }
    }
     \\ \\

    \theutterance \stepcounter{utterance}  
    & & & \multicolumn{2}{p{0.3\linewidth}}{
        \cellcolor[rgb]{0.9,0.9,0.9}{
            \makecell[{{p{\linewidth}}}]{
                \texttt{\tiny{[GM$|$GM]}}
                \texttt{continue} \\
            }
        }
    }
    & & \\ \\

    \theutterance \stepcounter{utterance}  
    & & & \multicolumn{2}{p{0.3\linewidth}}{
        \cellcolor[rgb]{0.9,0.9,0.9}{
            \makecell[{{p{\linewidth}}}]{
                \texttt{\tiny{[GM$|$GM]}}
                \texttt{<think>} \\
\texttt{Okay, let's start by understanding the problem. I need to divide the items between me and another player. The items are 2 hooks, 1 flower, 1 card, 1 eraser, 1 sponge. My value function is hook:3, flower:3, card:0, eraser:1, sponge:0. The other player's values are secret, so I don't know what they want. The goal is to maximize the total score for both of us, which means we need to find a division where our proposals don't conflict. If they do conflict, we both get zero, so it's crucial that our proposals are complementary.} \\
\\ 
\texttt{First, I should figure out what I value most. Hooks and flowers are the highest for me, each worth 3. The eraser is 1, and the others are 0. So ideally, I want to get as many hooks and flowers as possible. But there are only two hooks and one flower. If I take both hooks and the flower, that's 3+3+3=9 points for me. But then the other player would have to take the card, eraser, and sponge. However, maybe the other player values some of those items. For example, if they value the eraser or sponge, they might get some points from those. But if their values are different, maybe they want some of the items I have. But since I don't know their values, I need to communicate.} \\
\\ 
\texttt{The rules say we can exchange up to 5 messages each. So I should start by sending a message to the other player to discuss possible allocations. Let me think about what to propose. Since I don't know their preferences, maybe I can suggest a division that allows them to take items I don't value. For example, if I take the two hooks and the flower, I can offer the other player the eraser and sponge. But the card is also there. Alternatively, maybe I can take one hook and the flower, and let them take the other hook and the eraser. But I need to make sure that the other player can get their desired items without overlapping with mine.} \\
\\ 
\texttt{Wait, the key is to make sure that the proposals are complementary. So if I propose taking certain items, the other player's proposal must not require any of those items. For example, if I take both hooks and the flower, the other player can't take any of} \\
            }
        }
    }
    & & \\ \\

    \theutterance \stepcounter{utterance}  
    & & \multicolumn{4}{p{0.6\linewidth}}{
        \cellcolor[rgb]{0.9,0.9,0.9}{
            \makecell[{{p{\linewidth}}}]{
                \texttt{\tiny{[P1$\langle$GM]}}
                \texttt{<think>} \\
\texttt{Okay, let's start by understanding the problem. I need to divide the items between me and another player. The items are 2 hooks, 1 flower, 1 card, 1 eraser, 1 sponge. My value function is hook:3, flower:3, card:0, eraser:1, sponge:0. The other player's values are secret, so I don't know what they want. The goal is to maximize the total score for both of us, which means we need to find a division where our proposals don't conflict. If they do conflict, we both get zero, so it's crucial that our proposals are complementary.} \\
\\ 
\texttt{First, I should figure out what I value most. Hooks and flowers are the highest for me, each worth 3. The eraser is 1, and the others are 0. So ideally, I want to get as many hooks and flowers as possible. But there are only two hooks and one flower. If I take both hooks and the flower, that's 3+3+3=9 points for me. But then the other player would have to take the card, eraser, and sponge. However, maybe the other player values some of those items. For example, if they value the eraser or sponge, they might get some points from those. But if their values are different, maybe they want some of the items I have. But since I don't know their values, I need to communicate.} \\
\\ 
\texttt{The rules say we can exchange up to 5 messages each. So I should start by sending a message to the other player to discuss possible allocations. Let me think about what to propose. Since I don't know their preferences, maybe I can suggest a division that allows them to take items I don't value. For example, if I take the two hooks and the flower, I can offer the other player the eraser and sponge. But the card is also there. Alternatively, maybe I can take one hook and the flower, and let them take the other hook and the eraser. But I need to make sure that the other player can get their desired items without overlapping with mine.} \\
\\ 
\texttt{Wait, the key is to make sure that the proposals are complementary. So if I propose taking certain items, the other player's proposal must not require any of those items. For example, if I take both hooks and the flower, the other player can't take any of} \\
            }
        }
    }
    & \\ \\

    \theutterance \stepcounter{utterance}  
    & \multicolumn{4}{p{0.6\linewidth}}{
        \cellcolor[rgb]{0.9,0.9,0.9}{
            \makecell[{{p{\linewidth}}}]{
                \texttt{\tiny{[P1$\rangle$GM]}}
                \texttt{<think>} \\
\texttt{Okay, let's try to figure this out. I need to divide the items between me and the other player. The items are 2 hooks, 1 flower, 1 card, 1 eraser, 1 sponge. My values are hook:3, flower:3, eraser:1, others 0. The other player's values are secret, so I don't know what they want. The goal is to maximize the total score for both of us, so we need to find a way to split the items where our proposals don't conflict.} \\
\\ 
\texttt{First, I should consider what I value most. Hooks and flowers are the highest for me. So ideally, I want to get as many of those as possible. But there are only two hooks and one flower. If I take both hooks and the flower, that's 3+3+3=9 points for me. But then the other player would have to take the remaining items: card, eraser, sponge. But maybe they value some of those. However, if their values are different, maybe they want some of the items I have. But since I don't know their values, I need to communicate.} \\
\\ 
\texttt{The rules say we can exchange up to 5 messages each. So I should start by sending a message to the other player to discuss possible allocations. Let me think about what to propose. Since I don't know their preferences, maybe I can suggest a division that allows them to take items I don't value. For example, if I take the two hooks and the flower, I can offer the other player the eraser and sponge. But the card is also there. Alternatively, maybe I can take one hook and the flower, and let them take the other hook and the eraser. But I need to make sure that the other player can get their desired items without overlapping with mine.} \\
\\ 
\texttt{Wait, the key is to make sure that the proposals are complementary. So if I propose taking certain items, the other player's proposal must not require any of those items. For example, if I take both hooks and the flower, the other player can't take any of those. But if they want the eraser or sponge, that's okay. But if they also want a hook or the flower, that would conflict. So maybe I should propose taking one hook and the flower, and then offer the other hook and the eraser to them. That way, if they value the} \\
            }
        }
    }
    & & \\ \\

    \theutterance \stepcounter{utterance}  
    & & & \multicolumn{2}{p{0.3\linewidth}}{
        \cellcolor[rgb]{0.9,0.9,0.9}{
            \makecell[{{p{\linewidth}}}]{
                \texttt{\tiny{[GM$|$GM]}}
                \texttt{continue} \\
            }
        }
    }
    & & \\ \\

    \theutterance \stepcounter{utterance}  
    & & & \multicolumn{2}{p{0.3\linewidth}}{
        \cellcolor[rgb]{0.9,0.9,0.9}{
            \makecell[{{p{\linewidth}}}]{
                \texttt{\tiny{[GM$|$GM]}}
                \texttt{<think>} \\
\texttt{Okay, let's try to figure this out. I need to divide the items between me and the other player. The items are 2 hooks, 1 flower, 1 card, 1 eraser, 1 sponge. My values are hook:3, flower:3, eraser:1, others 0. The other player's values are secret, so I don't know what they want. The goal is to maximize the total score for both of us, so we need to find a way to split the items where our proposals don't conflict.} \\
\\ 
\texttt{First, I should consider what I value most. Hooks and flowers are the highest for me. So ideally, I want to get as many of those as possible. But there are only two hooks and one flower. If I take both hooks and the flower, that's 3+3+3=9 points for me. But then the other player would have to take the remaining items: card, eraser, sponge. But maybe they value some of those. However, if their values are different, maybe they want some of the items I have. But since I don't know their values, I need to communicate.} \\
\\ 
\texttt{The rules say we can exchange up to 5 messages each. So I should start by sending a message to the other player to discuss possible allocations. Let me think about what to propose. Since I don't know their preferences, maybe I can suggest a division that allows them to take items I don't value. For example, if I take the two hooks and the flower, I can offer the other player the eraser and sponge. But the card is also there. Alternatively, maybe I can take one hook and the flower, and let them take the other hook and the eraser. But I need to make sure that the other player can get their desired items without overlapping with mine.} \\
\\ 
\texttt{Wait, the key is to make sure that the proposals are complementary. So if I propose taking certain items, the other player's proposal must not require any of those items. For example, if I take both hooks and the flower, the other player can't take any of those. But if they want the eraser or sponge, that's okay. But if they also want a hook or the flower, that would conflict. So maybe I should propose taking one hook and the flower, and then offer the other hook and the eraser to them. That way, if they value the} \\
            }
        }
    }
    & & \\ \\

    \theutterance \stepcounter{utterance}  
    & & \multicolumn{4}{p{0.6\linewidth}}{
        \cellcolor[rgb]{0.9,0.9,0.9}{
            \makecell[{{p{\linewidth}}}]{
                \texttt{\tiny{[GM$\rangle$P2]}}
                \texttt{<think>} \\
\texttt{Okay, let's try to figure this out. I need to divide the items between me and the other player. The items are 2 hooks, 1 flower, 1 card, 1 eraser, 1 sponge. My values are hook:3, flower:3, eraser:1, others 0. The other player's values are secret, so I don't know what they want. The goal is to maximize the total score for both of us, so we need to find a way to split the items where our proposals don't conflict.} \\
\\ 
\texttt{First, I should consider what I value most. Hooks and flowers are the highest for me. So ideally, I want to get as many of those as possible. But there are only two hooks and one flower. If I take both hooks and the flower, that's 3+3+3=9 points for me. But then the other player would have to take the remaining items: card, eraser, sponge. But maybe they value some of those. However, if their values are different, maybe they want some of the items I have. But since I don't know their values, I need to communicate.} \\
\\ 
\texttt{The rules say we can exchange up to 5 messages each. So I should start by sending a message to the other player to discuss possible allocations. Let me think about what to propose. Since I don't know their preferences, maybe I can suggest a division that allows them to take items I don't value. For example, if I take the two hooks and the flower, I can offer the other player the eraser and sponge. But the card is also there. Alternatively, maybe I can take one hook and the flower, and let them take the other hook and the eraser. But I need to make sure that the other player can get their desired items without overlapping with mine.} \\
\\ 
\texttt{Wait, the key is to make sure that the proposals are complementary. So if I propose taking certain items, the other player's proposal must not require any of those items. For example, if I take both hooks and the flower, the other player can't take any of those. But if they want the eraser or sponge, that's okay. But if they also want a hook or the flower, that would conflict. So maybe I should propose taking one hook and the flower, and then offer the other hook and the eraser to them. That way, if they value the} \\
            }
        }
    }
    & \\ \\

    \theutterance \stepcounter{utterance}  
    & & & \multicolumn{4}{p{0.6\linewidth}}{
        \cellcolor[rgb]{0.9,0.9,0.9}{
            \makecell[{{p{\linewidth}}}]{
                \texttt{\tiny{[GM$\langle$P2]}}
                \texttt{<think>} \\
\texttt{Okay, let's try to figure this out. I need to divide the items between me and the other player. The items are 2 hooks, 1 flower, 1 card, 1 eraser, 1 sponge. My values are hook:3, flower:3, eraser:1, others 0. The other player's values are secret, so I don't know what they want. The goal is to maximize the total score for both of us, so we need to find a way to split the items where our proposals don't conflict.} \\
\\ 
\texttt{First, I should consider what I value most. Hooks and flowers are the highest for me. So ideally, I want to get as many of those as possible. But there are only two hooks and one flower. If I take both hooks and the flower, that's 3+3+3=9 points for me. But then the other player would have to take the remaining items: card, eraser, sponge. But maybe they value some of those. However, if their values are different, maybe they want some of the items I have. But since I don't know their values, I need to communicate.} \\
\\ 
\texttt{The rules say we can exchange up to 5 messages each. So I should start by sending a message to the other player to discuss possible allocations. Let me think about what to propose. Since I don't know their preferences, maybe I can suggest a division that allows them to take items I don't value. For example, if I take the two hooks and the flower, I can offer the other player the eraser and sponge. But the card is also there. Alternatively, maybe I can take one hook and the flower, and let them take the other hook and the eraser. But I need to make sure that the other player can get their desired items without overlapping with mine.} \\
\\ 
\texttt{Wait, the key is to make sure that the proposals are complementary. So if I propose taking certain items, the other player's proposal must not require any of those items. For example, if I take both hooks and the flower, the other player can't take any of those. But if they want the eraser or sponge, that's okay. But if they also want a hook or the flower, that would conflict. So maybe I should propose taking one hook and the flower, and then offer the other hook and the eraser to them. That way, if they value the} \\
            }
        }
    }
     \\ \\

    \theutterance \stepcounter{utterance}  
    & & & \multicolumn{2}{p{0.3\linewidth}}{
        \cellcolor[rgb]{0.9,0.9,0.9}{
            \makecell[{{p{\linewidth}}}]{
                \texttt{\tiny{[GM$|$GM]}}
                \texttt{continue} \\
            }
        }
    }
    & & \\ \\

    \theutterance \stepcounter{utterance}  
    & & & \multicolumn{2}{p{0.3\linewidth}}{
        \cellcolor[rgb]{0.9,0.9,0.9}{
            \makecell[{{p{\linewidth}}}]{
                \texttt{\tiny{[GM$|$GM]}}
                \texttt{<think>} \\
\texttt{Okay, let's try to figure this out. I need to divide the items between me and the other player. The items are 2 hooks, 1 flower, 1 card, 1 eraser, 1 sponge. My values are hook:3, flower:3, eraser:1, others 0. The other player's values are secret, so I don't know what they want. The goal is to maximize the total score for both of us, so we need to find a way to split the items where our proposals don't conflict.} \\
\\ 
\texttt{First, I should consider what I value most. Hooks and flowers are the highest for me. So ideally, I want to get as many of those as possible. But there are only two hooks and one flower. If I take both hooks and the flower, that's 3+3+3=9 points for me. But then the other player would have to take the remaining items: card, eraser, sponge. But maybe they value some of those. However, if their values are different, maybe they want some of the items I have. But since I don't know their values, I need to communicate.} \\
\\ 
\texttt{The rules say we can exchange up to 5 messages each. So I should start by sending a message to the other player to discuss possible allocations. Let me think about what to propose. Since I don't know their preferences, maybe I can suggest a division that allows them to take items I don't value. For example, if I take the two hooks and the flower, I can offer the other player the eraser and sponge. But the card is also there. Alternatively, maybe I can take one hook and the flower, and let them take the other hook and the eraser. But I need to make sure that the other player can get their desired items without overlapping with mine.} \\
\\ 
\texttt{Wait, the key is to make sure that the proposals are complementary. So if I propose taking certain items, the other player's proposal must not require any of those items. For example, if I take both hooks and the flower, the other player can't take any of those. But if they want the eraser or sponge, that's okay. But if they also want a hook or the flower, that would conflict. So maybe I should propose taking one hook and the flower, and then offer the other hook and the eraser to them. That way, if they value the} \\
            }
        }
    }
    & & \\ \\

    \theutterance \stepcounter{utterance}  
    & & \multicolumn{4}{p{0.6\linewidth}}{
        \cellcolor[rgb]{0.9,0.9,0.9}{
            \makecell[{{p{\linewidth}}}]{
                \texttt{\tiny{[P1$\langle$GM]}}
                \texttt{<think>} \\
\texttt{Okay, let's try to figure this out. I need to divide the items between me and the other player. The items are 2 hooks, 1 flower, 1 card, 1 eraser, 1 sponge. My values are hook:3, flower:3, eraser:1, others 0. The other player's values are secret, so I don't know what they want. The goal is to maximize the total score for both of us, so we need to find a way to split the items where our proposals don't conflict.} \\
\\ 
\texttt{First, I should consider what I value most. Hooks and flowers are the highest for me. So ideally, I want to get as many of those as possible. But there are only two hooks and one flower. If I take both hooks and the flower, that's 3+3+3=9 points for me. But then the other player would have to take the remaining items: card, eraser, sponge. But maybe they value some of those. However, if their values are different, maybe they want some of the items I have. But since I don't know their values, I need to communicate.} \\
\\ 
\texttt{The rules say we can exchange up to 5 messages each. So I should start by sending a message to the other player to discuss possible allocations. Let me think about what to propose. Since I don't know their preferences, maybe I can suggest a division that allows them to take items I don't value. For example, if I take the two hooks and the flower, I can offer the other player the eraser and sponge. But the card is also there. Alternatively, maybe I can take one hook and the flower, and let them take the other hook and the eraser. But I need to make sure that the other player can get their desired items without overlapping with mine.} \\
\\ 
\texttt{Wait, the key is to make sure that the proposals are complementary. So if I propose taking certain items, the other player's proposal must not require any of those items. For example, if I take both hooks and the flower, the other player can't take any of those. But if they want the eraser or sponge, that's okay. But if they also want a hook or the flower, that would conflict. So maybe I should propose taking one hook and the flower, and then offer the other hook and the eraser to them. That way, if they value the} \\
            }
        }
    }
    & \\ \\

    \theutterance \stepcounter{utterance}  
    & \multicolumn{4}{p{0.6\linewidth}}{
        \cellcolor[rgb]{0.9,0.9,0.9}{
            \makecell[{{p{\linewidth}}}]{
                \texttt{\tiny{[P1$\rangle$GM]}}
                \texttt{<think>} \\
\texttt{Okay, let's try to figure this out. I need to divide the items between me and the other player. The items are 2 hooks, 1 flower, 1 card, 1 eraser, 1 sponge. My values are hook:3, flower:3, eraser:1, others 0. The other player's values are secret, so I don't know what they want. The goal is to maximize the total score for both of us, so we need to find a way to split the items where our proposals don't conflict.} \\
\\ 
\texttt{First, I should consider what I value most. Hooks and flowers are the highest for me. So ideally, I want to get as many of those as possible. But there are only two hooks and one flower. If I take both hooks and the flower, that's 3+3+3=9 points for me. But then the other player would have to take the remaining items: card, eraser, sponge. But maybe they value some of those. However, if their values are different, maybe they want some of the items I have. But since I don't know their values, I need to communicate.} \\
\\ 
\texttt{The rules say we can exchange up to 5 messages each. So I should start by sending a message to the other player to discuss possible allocations. Let me think about what to propose. Since I don't know their preferences, maybe I can suggest a division that allows them to take items I don't value. For example, if I take the two hooks and the flower, I can offer the other player the eraser and sponge. But the card is also there. Alternatively, maybe I can take one hook and the flower, and let them take the other hook and the eraser. But I need to make sure that the other player can get their desired items without overlapping with mine.} \\
\\ 
\texttt{Wait, the key is to make sure that the proposals are complementary. So if I propose taking certain items, the other player's proposal must not require any of those items. For example, if I take both hooks and the flower, the other player can't take any of those. But if they want the eraser or sponge, that's okay. But if they also want a hook or the flower, that would conflict. So maybe I should propose taking one hook and the flower, and then offer the other hook and the eraser to them. That way, if they value the} \\
            }
        }
    }
    & & \\ \\

    \theutterance \stepcounter{utterance}  
    & & & \multicolumn{2}{p{0.3\linewidth}}{
        \cellcolor[rgb]{0.9,0.9,0.9}{
            \makecell[{{p{\linewidth}}}]{
                \texttt{\tiny{[GM$|$GM]}}
                \texttt{continue} \\
            }
        }
    }
    & & \\ \\

    \theutterance \stepcounter{utterance}  
    & & & \multicolumn{2}{p{0.3\linewidth}}{
        \cellcolor[rgb]{0.9,0.9,0.9}{
            \makecell[{{p{\linewidth}}}]{
                \texttt{\tiny{[GM$|$GM]}}
                \texttt{<think>} \\
\texttt{Okay, let's try to figure this out. I need to divide the items between me and the other player. The items are 2 hooks, 1 flower, 1 card, 1 eraser, 1 sponge. My values are hook:3, flower:3, eraser:1, others 0. The other player's values are secret, so I don't know what they want. The goal is to maximize the total score for both of us, so we need to find a way to split the items where our proposals don't conflict.} \\
\\ 
\texttt{First, I should consider what I value most. Hooks and flowers are the highest for me. So ideally, I want to get as many of those as possible. But there are only two hooks and one flower. If I take both hooks and the flower, that's 3+3+3=9 points for me. But then the other player would have to take the remaining items: card, eraser, sponge. But maybe they value some of those. However, if their values are different, maybe they want some of the items I have. But since I don't know their values, I need to communicate.} \\
\\ 
\texttt{The rules say we can exchange up to 5 messages each. So I should start by sending a message to the other player to discuss possible allocations. Let me think about what to propose. Since I don't know their preferences, maybe I can suggest a division that allows them to take items I don't value. For example, if I take the two hooks and the flower, I can offer the other player the eraser and sponge. But the card is also there. Alternatively, maybe I can take one hook and the flower, and let them take the other hook and the eraser. But I need to make sure that the other player can get their desired items without overlapping with mine.} \\
\\ 
\texttt{Wait, the key is to make sure that the proposals are complementary. So if I propose taking certain items, the other player's proposal must not require any of those items. For example, if I take both hooks and the flower, the other player can't take any of those. But if they want the eraser or sponge, that's okay. But if they also want a hook or the flower, that would conflict. So maybe I should propose taking one hook and the flower, and then offer the other hook and the eraser to them. That way, if they value the} \\
            }
        }
    }
    & & \\ \\

    \theutterance \stepcounter{utterance}  
    & & \multicolumn{4}{p{0.6\linewidth}}{
        \cellcolor[rgb]{0.9,0.9,0.9}{
            \makecell[{{p{\linewidth}}}]{
                \texttt{\tiny{[GM$\rangle$P2]}}
                \texttt{<think>} \\
\texttt{Okay, let's try to figure this out. I need to divide the items between me and the other player. The items are 2 hooks, 1 flower, 1 card, 1 eraser, 1 sponge. My values are hook:3, flower:3, eraser:1, others 0. The other player's values are secret, so I don't know what they want. The goal is to maximize the total score for both of us, so we need to find a way to split the items where our proposals don't conflict.} \\
\\ 
\texttt{First, I should consider what I value most. Hooks and flowers are the highest for me. So ideally, I want to get as many of those as possible. But there are only two hooks and one flower. If I take both hooks and the flower, that's 3+3+3=9 points for me. But then the other player would have to take the remaining items: card, eraser, sponge. But maybe they value some of those. However, if their values are different, maybe they want some of the items I have. But since I don't know their values, I need to communicate.} \\
\\ 
\texttt{The rules say we can exchange up to 5 messages each. So I should start by sending a message to the other player to discuss possible allocations. Let me think about what to propose. Since I don't know their preferences, maybe I can suggest a division that allows them to take items I don't value. For example, if I take the two hooks and the flower, I can offer the other player the eraser and sponge. But the card is also there. Alternatively, maybe I can take one hook and the flower, and let them take the other hook and the eraser. But I need to make sure that the other player can get their desired items without overlapping with mine.} \\
\\ 
\texttt{Wait, the key is to make sure that the proposals are complementary. So if I propose taking certain items, the other player's proposal must not require any of those items. For example, if I take both hooks and the flower, the other player can't take any of those. But if they want the eraser or sponge, that's okay. But if they also want a hook or the flower, that would conflict. So maybe I should propose taking one hook and the flower, and then offer the other hook and the eraser to them. That way, if they value the} \\
            }
        }
    }
    & \\ \\

    \theutterance \stepcounter{utterance}  
    & & & \multicolumn{4}{p{0.6\linewidth}}{
        \cellcolor[rgb]{0.9,0.9,0.9}{
            \makecell[{{p{\linewidth}}}]{
                \texttt{\tiny{[GM$\langle$P2]}}
                \texttt{<think>} \\
\texttt{Okay, let's try to figure this out. I need to divide the items between me and the other player. The items are 2 hooks, 1 flower, 1 card, 1 eraser, 1 sponge. My values are hook:3, flower:3, eraser:1, others 0. The other player's values are secret, so I don't know what they want. The goal is to maximize the total score for both of us, so we need to find a way to split the items where our proposals don't conflict.} \\
\\ 
\texttt{First, I should consider what I value most. Hooks and flowers are the highest for me. So ideally, I want to get as many of those as possible. But there are only two hooks and one flower. If I take both hooks and the flower, that's 3+3+3=9 points for me. But then the other player would have to take the remaining items: card, eraser, sponge. But maybe they value some of those. However, if their values are different, maybe they want some of the items I have. But since I don't know their values, I need to communicate.} \\
\\ 
\texttt{The rules say we can exchange up to 5 messages each. So I should start by sending a message to the other player to discuss possible allocations. Let me think about what to propose. Since I don't know their preferences, maybe I can suggest a division that allows them to take items I don't value. For example, if I take the two hooks and the flower, I can offer the other player the eraser and sponge. But the card is also there. Alternatively, maybe I can take one hook and the flower, and let them take the other hook and the eraser. But I need to make sure that the other player can get their desired items without overlapping with mine.} \\
\\ 
\texttt{Wait, the key is to make sure that the proposals are complementary. So if I propose taking certain items, the other player's proposal must not require any of those items. For example, if I take both hooks and the flower, the other player can't take any of those. But if they want the eraser or sponge, that's okay. But if they also want a hook or the flower, that would conflict. So maybe I should propose taking one hook and the flower, and then offer the other hook and the eraser to them. That way, if they value the} \\
            }
        }
    }
     \\ \\

    \theutterance \stepcounter{utterance}  
    & & & \multicolumn{2}{p{0.3\linewidth}}{
        \cellcolor[rgb]{0.9,0.9,0.9}{
            \makecell[{{p{\linewidth}}}]{
                \texttt{\tiny{[GM$|$GM]}}
                \texttt{continue} \\
            }
        }
    }
    & & \\ \\

    \theutterance \stepcounter{utterance}  
    & & & \multicolumn{2}{p{0.3\linewidth}}{
        \cellcolor[rgb]{0.9,0.9,0.9}{
            \makecell[{{p{\linewidth}}}]{
                \texttt{\tiny{[GM$|$GM]}}
                \texttt{<think>} \\
\texttt{Okay, let's try to figure this out. I need to divide the items between me and the other player. The items are 2 hooks, 1 flower, 1 card, 1 eraser, 1 sponge. My values are hook:3, flower:3, eraser:1, others 0. The other player's values are secret, so I don't know what they want. The goal is to maximize the total score for both of us, so we need to find a way to split the items where our proposals don't conflict.} \\
\\ 
\texttt{First, I should consider what I value most. Hooks and flowers are the highest for me. So ideally, I want to get as many of those as possible. But there are only two hooks and one flower. If I take both hooks and the flower, that's 3+3+3=9 points for me. But then the other player would have to take the remaining items: card, eraser, sponge. But maybe they value some of those. However, if their values are different, maybe they want some of the items I have. But since I don't know their values, I need to communicate.} \\
\\ 
\texttt{The rules say we can exchange up to 5 messages each. So I should start by sending a message to the other player to discuss possible allocations. Let me think about what to propose. Since I don't know their preferences, maybe I can suggest a division that allows them to take items I don't value. For example, if I take the two hooks and the flower, I can offer the other player the eraser and sponge. But the card is also there. Alternatively, maybe I can take one hook and the flower, and let them take the other hook and the eraser. But I need to make sure that the other player can get their desired items without overlapping with mine.} \\
\\ 
\texttt{Wait, the key is to make sure that the proposals are complementary. So if I propose taking certain items, the other player's proposal must not require any of those items. For example, if I take both hooks and the flower, the other player can't take any of those. But if they want the eraser or sponge, that's okay. But if they also want a hook or the flower, that would conflict. So maybe I should propose taking one hook and the flower, and then offer the other hook and the eraser to them. That way, if they value the} \\
            }
        }
    }
    & & \\ \\

    \theutterance \stepcounter{utterance}  
    & & \multicolumn{4}{p{0.6\linewidth}}{
        \cellcolor[rgb]{0.9,0.9,0.9}{
            \makecell[{{p{\linewidth}}}]{
                \texttt{\tiny{[P1$\langle$GM]}}
                \texttt{<think>} \\
\texttt{Okay, let's try to figure this out. I need to divide the items between me and the other player. The items are 2 hooks, 1 flower, 1 card, 1 eraser, 1 sponge. My values are hook:3, flower:3, eraser:1, others 0. The other player's values are secret, so I don't know what they want. The goal is to maximize the total score for both of us, so we need to find a way to split the items where our proposals don't conflict.} \\
\\ 
\texttt{First, I should consider what I value most. Hooks and flowers are the highest for me. So ideally, I want to get as many of those as possible. But there are only two hooks and one flower. If I take both hooks and the flower, that's 3+3+3=9 points for me. But then the other player would have to take the remaining items: card, eraser, sponge. But maybe they value some of those. However, if their values are different, maybe they want some of the items I have. But since I don't know their values, I need to communicate.} \\
\\ 
\texttt{The rules say we can exchange up to 5 messages each. So I should start by sending a message to the other player to discuss possible allocations. Let me think about what to propose. Since I don't know their preferences, maybe I can suggest a division that allows them to take items I don't value. For example, if I take the two hooks and the flower, I can offer the other player the eraser and sponge. But the card is also there. Alternatively, maybe I can take one hook and the flower, and let them take the other hook and the eraser. But I need to make sure that the other player can get their desired items without overlapping with mine.} \\
\\ 
\texttt{Wait, the key is to make sure that the proposals are complementary. So if I propose taking certain items, the other player's proposal must not require any of those items. For example, if I take both hooks and the flower, the other player can't take any of those. But if they want the eraser or sponge, that's okay. But if they also want a hook or the flower, that would conflict. So maybe I should propose taking one hook and the flower, and then offer the other hook and the eraser to them. That way, if they value the} \\
            }
        }
    }
    & \\ \\

    \theutterance \stepcounter{utterance}  
    & \multicolumn{4}{p{0.6\linewidth}}{
        \cellcolor[rgb]{0.9,0.9,0.9}{
            \makecell[{{p{\linewidth}}}]{
                \texttt{\tiny{[P1$\rangle$GM]}}
                \texttt{<think>} \\
\texttt{Okay, let's try to figure this out. I need to divide the items between me and the other player. The items are 2 hooks, 1 flower, 1 card, 1 eraser, 1 sponge. My values are hook:3, flower:3, eraser:1, others 0. The other player's values are secret, so I don't know what they want. The goal is to maximize the total score for both of us, so we need to find a way to split the items where our proposals don't conflict.} \\
\\ 
\texttt{First, I should consider what I value most. Hooks and flowers are the highest for me. So ideally, I want to get as many of those as possible. But there are only two hooks and one flower. If I take both hooks and the flower, that's 3+3+3=9 points for me. But then the other player would have to take the remaining items: card, eraser, sponge. But maybe they value some of those. However, if their values are different, maybe they want some of the items I have. But since I don't know their values, I need to communicate.} \\
\\ 
\texttt{The rules say we can exchange up to 5 messages each. So I should start by sending a message to the other player to discuss possible allocations. Let me think about what to propose. Since I don't know their preferences, maybe I can suggest a division that allows them to take items I don't value. For example, if I take the two hooks and the flower, I can offer the other player the eraser and sponge. But the card is also there. Alternatively, maybe I can take one hook and the flower, and let them take the other hook and the eraser. But I need to make sure that the other player can get their desired items without overlapping with mine.} \\
\\ 
\texttt{Wait, the key is to make sure that the proposals are complementary. So if I propose taking certain items, the other player's proposal must not require any of those items. For example, if I take both hooks and the flower, the other player can't take any of those. But if they want the eraser or sponge, that's okay. But if they also want a hook or the flower, that would conflict. So maybe I should propose taking one hook and the flower, and then offer the other hook and the eraser to them. That way, if they value the} \\
            }
        }
    }
    & & \\ \\

    \theutterance \stepcounter{utterance}  
    & & & \multicolumn{2}{p{0.3\linewidth}}{
        \cellcolor[rgb]{0.9,0.9,0.9}{
            \makecell[{{p{\linewidth}}}]{
                \texttt{\tiny{[GM$|$GM]}}
                \texttt{continue} \\
            }
        }
    }
    & & \\ \\

    \theutterance \stepcounter{utterance}  
    & & & \multicolumn{2}{p{0.3\linewidth}}{
        \cellcolor[rgb]{0.9,0.9,0.9}{
            \makecell[{{p{\linewidth}}}]{
                \texttt{\tiny{[GM$|$GM]}}
                \texttt{<think>} \\
\texttt{Okay, let's try to figure this out. I need to divide the items between me and the other player. The items are 2 hooks, 1 flower, 1 card, 1 eraser, 1 sponge. My values are hook:3, flower:3, eraser:1, others 0. The other player's values are secret, so I don't know what they want. The goal is to maximize the total score for both of us, so we need to find a way to split the items where our proposals don't conflict.} \\
\\ 
\texttt{First, I should consider what I value most. Hooks and flowers are the highest for me. So ideally, I want to get as many of those as possible. But there are only two hooks and one flower. If I take both hooks and the flower, that's 3+3+3=9 points for me. But then the other player would have to take the remaining items: card, eraser, sponge. But maybe they value some of those. However, if their values are different, maybe they want some of the items I have. But since I don't know their values, I need to communicate.} \\
\\ 
\texttt{The rules say we can exchange up to 5 messages each. So I should start by sending a message to the other player to discuss possible allocations. Let me think about what to propose. Since I don't know their preferences, maybe I can suggest a division that allows them to take items I don't value. For example, if I take the two hooks and the flower, I can offer the other player the eraser and sponge. But the card is also there. Alternatively, maybe I can take one hook and the flower, and let them take the other hook and the eraser. But I need to make sure that the other player can get their desired items without overlapping with mine.} \\
\\ 
\texttt{Wait, the key is to make sure that the proposals are complementary. So if I propose taking certain items, the other player's proposal must not require any of those items. For example, if I take both hooks and the flower, the other player can't take any of those. But if they want the eraser or sponge, that's okay. But if they also want a hook or the flower, that would conflict. So maybe I should propose taking one hook and the flower, and then offer the other hook and the eraser to them. That way, if they value the} \\
            }
        }
    }
    & & \\ \\

    \theutterance \stepcounter{utterance}  
    & & \multicolumn{4}{p{0.6\linewidth}}{
        \cellcolor[rgb]{0.9,0.9,0.9}{
            \makecell[{{p{\linewidth}}}]{
                \texttt{\tiny{[GM$\rangle$P2]}}
                \texttt{<think>} \\
\texttt{Okay, let's try to figure this out. I need to divide the items between me and the other player. The items are 2 hooks, 1 flower, 1 card, 1 eraser, 1 sponge. My values are hook:3, flower:3, eraser:1, others 0. The other player's values are secret, so I don't know what they want. The goal is to maximize the total score for both of us, so we need to find a way to split the items where our proposals don't conflict.} \\
\\ 
\texttt{First, I should consider what I value most. Hooks and flowers are the highest for me. So ideally, I want to get as many of those as possible. But there are only two hooks and one flower. If I take both hooks and the flower, that's 3+3+3=9 points for me. But then the other player would have to take the remaining items: card, eraser, sponge. But maybe they value some of those. However, if their values are different, maybe they want some of the items I have. But since I don't know their values, I need to communicate.} \\
\\ 
\texttt{The rules say we can exchange up to 5 messages each. So I should start by sending a message to the other player to discuss possible allocations. Let me think about what to propose. Since I don't know their preferences, maybe I can suggest a division that allows them to take items I don't value. For example, if I take the two hooks and the flower, I can offer the other player the eraser and sponge. But the card is also there. Alternatively, maybe I can take one hook and the flower, and let them take the other hook and the eraser. But I need to make sure that the other player can get their desired items without overlapping with mine.} \\
\\ 
\texttt{Wait, the key is to make sure that the proposals are complementary. So if I propose taking certain items, the other player's proposal must not require any of those items. For example, if I take both hooks and the flower, the other player can't take any of those. But if they want the eraser or sponge, that's okay. But if they also want a hook or the flower, that would conflict. So maybe I should propose taking one hook and the flower, and then offer the other hook and the eraser to them. That way, if they value the} \\
            }
        }
    }
    & \\ \\

    \theutterance \stepcounter{utterance}  
    & & & \multicolumn{4}{p{0.6\linewidth}}{
        \cellcolor[rgb]{0.9,0.9,0.9}{
            \makecell[{{p{\linewidth}}}]{
                \texttt{\tiny{[GM$\langle$P2]}}
                \texttt{<think>} \\
\texttt{Okay, let's try to figure this out. I need to divide the items between me and the other player. The items are 2 hooks, 1 flower, 1 card, 1 eraser, 1 sponge. My values are hook:3, flower:3, eraser:1, others 0. The other player's values are secret, so I don't know what they want. The goal is to maximize the total score for both of us, so we need to find a way to split the items where our proposals don't conflict.} \\
\\ 
\texttt{First, I should consider what I value most. Hooks and flowers are the highest for me. So ideally, I want to get as many of those as possible. But there are only two hooks and one flower. If I take both hooks and the flower, that's 3+3+3=9 points for me. But then the other player would have to take the remaining items: card, eraser, sponge. But maybe they value some of those. However, if their values are different, maybe they want some of the items I have. But since I don't know their values, I need to communicate.} \\
\\ 
\texttt{The rules say we can exchange up to 5 messages each. So I should start by sending a message to the other player to discuss possible allocations. Let me think about what to propose. Since I don't know their preferences, maybe I can suggest a division that allows them to take items I don't value. For example, if I take the two hooks and the flower, I can offer the other player the eraser and sponge. But the card is also there. Alternatively, maybe I can take one hook and the flower, and let them take the other hook and the eraser. But I need to make sure that the other player can get their desired items without overlapping with mine.} \\
\\ 
\texttt{Wait, the key is to make sure that the proposals are complementary. So if I propose taking certain items, the other player's proposal must not require any of those items. For example, if I take both hooks and the flower, the other player can't take any of those. But if they want the eraser or sponge, that's okay. But if they also want a hook or the flower, that would conflict. So maybe I should propose taking one hook and the flower, and then offer the other hook and the eraser to them. That way, if they value the} \\
            }
        }
    }
     \\ \\

    \theutterance \stepcounter{utterance}  
    & & & \multicolumn{2}{p{0.3\linewidth}}{
        \cellcolor[rgb]{0.9,0.9,0.9}{
            \makecell[{{p{\linewidth}}}]{
                \texttt{\tiny{[GM$|$GM]}}
                \texttt{continue} \\
            }
        }
    }
    & & \\ \\

    \theutterance \stepcounter{utterance}  
    & & & \multicolumn{2}{p{0.3\linewidth}}{
        \cellcolor[rgb]{0.9,0.9,0.9}{
            \makecell[{{p{\linewidth}}}]{
                \texttt{\tiny{[GM$|$GM]}}
                \texttt{<think>} \\
\texttt{Okay, let's try to figure this out. I need to divide the items between me and the other player. The items are 2 hooks, 1 flower, 1 card, 1 eraser, 1 sponge. My values are hook:3, flower:3, eraser:1, others 0. The other player's values are secret, so I don't know what they want. The goal is to maximize the total score for both of us, so we need to find a way to split the items where our proposals don't conflict.} \\
\\ 
\texttt{First, I should consider what I value most. Hooks and flowers are the highest for me. So ideally, I want to get as many of those as possible. But there are only two hooks and one flower. If I take both hooks and the flower, that's 3+3+3=9 points for me. But then the other player would have to take the remaining items: card, eraser, sponge. But maybe they value some of those. However, if their values are different, maybe they want some of the items I have. But since I don't know their values, I need to communicate.} \\
\\ 
\texttt{The rules say we can exchange up to 5 messages each. So I should start by sending a message to the other player to discuss possible allocations. Let me think about what to propose. Since I don't know their preferences, maybe I can suggest a division that allows them to take items I don't value. For example, if I take the two hooks and the flower, I can offer the other player the eraser and sponge. But the card is also there. Alternatively, maybe I can take one hook and the flower, and let them take the other hook and the eraser. But I need to make sure that the other player can get their desired items without overlapping with mine.} \\
\\ 
\texttt{Wait, the key is to make sure that the proposals are complementary. So if I propose taking certain items, the other player's proposal must not require any of those items. For example, if I take both hooks and the flower, the other player can't take any of those. But if they want the eraser or sponge, that's okay. But if they also want a hook or the flower, that would conflict. So maybe I should propose taking one hook and the flower, and then offer the other hook and the eraser to them. That way, if they value the} \\
            }
        }
    }
    & & \\ \\

    \theutterance \stepcounter{utterance}  
    & & \multicolumn{4}{p{0.6\linewidth}}{
        \cellcolor[rgb]{0.9,0.9,0.9}{
            \makecell[{{p{\linewidth}}}]{
                \texttt{\tiny{[P1$\langle$GM]}}
                \texttt{<think>} \\
\texttt{Okay, let's try to figure this out. I need to divide the items between me and the other player. The items are 2 hooks, 1 flower, 1 card, 1 eraser, 1 sponge. My values are hook:3, flower:3, eraser:1, others 0. The other player's values are secret, so I don't know what they want. The goal is to maximize the total score for both of us, so we need to find a way to split the items where our proposals don't conflict.} \\
\\ 
\texttt{First, I should consider what I value most. Hooks and flowers are the highest for me. So ideally, I want to get as many of those as possible. But there are only two hooks and one flower. If I take both hooks and the flower, that's 3+3+3=9 points for me. But then the other player would have to take the remaining items: card, eraser, sponge. But maybe they value some of those. However, if their values are different, maybe they want some of the items I have. But since I don't know their values, I need to communicate.} \\
\\ 
\texttt{The rules say we can exchange up to 5 messages each. So I should start by sending a message to the other player to discuss possible allocations. Let me think about what to propose. Since I don't know their preferences, maybe I can suggest a division that allows them to take items I don't value. For example, if I take the two hooks and the flower, I can offer the other player the eraser and sponge. But the card is also there. Alternatively, maybe I can take one hook and the flower, and let them take the other hook and the eraser. But I need to make sure that the other player can get their desired items without overlapping with mine.} \\
\\ 
\texttt{Wait, the key is to make sure that the proposals are complementary. So if I propose taking certain items, the other player's proposal must not require any of those items. For example, if I take both hooks and the flower, the other player can't take any of those. But if they want the eraser or sponge, that's okay. But if they also want a hook or the flower, that would conflict. So maybe I should propose taking one hook and the flower, and then offer the other hook and the eraser to them. That way, if they value the} \\
            }
        }
    }
    & \\ \\

    \theutterance \stepcounter{utterance}  
    & \multicolumn{4}{p{0.6\linewidth}}{
        \cellcolor[rgb]{0.9,0.9,0.9}{
            \makecell[{{p{\linewidth}}}]{
                \texttt{\tiny{[P1$\rangle$GM]}}
                \texttt{<think>} \\
\texttt{Okay, let's try to figure this out. I need to divide the items between me and the other player. The items are 2 hooks, 1 flower, 1 card, 1 eraser, 1 sponge. My values are hook:3, flower:3, eraser:1, others 0. The other player's values are secret, so I don't know what they want. The goal is to maximize the total score for both of us, so we need to find a way to split the items where our proposals don't conflict.} \\
\\ 
\texttt{First, I should consider what I value most. Hooks and flowers are the highest for me. So ideally, I want to get as many of those as possible. But there are only two hooks and one flower. If I take both hooks and the flower, that's 3+3+3=9 points for me. But then the other player would have to take the remaining items: card, eraser, sponge. But maybe they value some of those. However, if their values are different, maybe they want some of the items I have. But since I don't know their values, I need to communicate.} \\
\\ 
\texttt{The rules say we can exchange up to 5 messages each. So I should start by sending a message to the other player to discuss possible allocations. Let me think about what to propose. Since I don't know their preferences, maybe I can suggest a division that allows them to take items I don't value. For example, if I take the two hooks and the flower, I can offer the other player the eraser and sponge. But the card is also there. Alternatively, maybe I can take one hook and the flower, and let them take the other hook and the eraser. But I need to make sure that the other player can get their desired items without overlapping with mine.} \\
\\ 
\texttt{Wait, the key is to make sure that the proposals are complementary. So if I propose taking certain items, the other player's proposal must not require any of those items. For example, if I take both hooks and the flower, the other player can't take any of those. But if they want the eraser or sponge, that's okay. But if they also want a hook or the flower, that would conflict. So maybe I should propose taking one hook and the flower, and then offer the other hook and the eraser to them. That way, if they value the} \\
            }
        }
    }
    & & \\ \\

    \theutterance \stepcounter{utterance}  
    & & & \multicolumn{2}{p{0.3\linewidth}}{
        \cellcolor[rgb]{0.9,0.9,0.9}{
            \makecell[{{p{\linewidth}}}]{
                \texttt{\tiny{[GM$|$GM]}}
                \texttt{continue} \\
            }
        }
    }
    & & \\ \\

    \theutterance \stepcounter{utterance}  
    & & & \multicolumn{2}{p{0.3\linewidth}}{
        \cellcolor[rgb]{0.9,0.9,0.9}{
            \makecell[{{p{\linewidth}}}]{
                \texttt{\tiny{[GM$|$GM]}}
                \texttt{<think>} \\
\texttt{Okay, let's try to figure this out. I need to divide the items between me and the other player. The items are 2 hooks, 1 flower, 1 card, 1 eraser, 1 sponge. My values are hook:3, flower:3, eraser:1, others 0. The other player's values are secret, so I don't know what they want. The goal is to maximize the total score for both of us, so we need to find a way to split the items where our proposals don't conflict.} \\
\\ 
\texttt{First, I should consider what I value most. Hooks and flowers are the highest for me. So ideally, I want to get as many of those as possible. But there are only two hooks and one flower. If I take both hooks and the flower, that's 3+3+3=9 points for me. But then the other player would have to take the remaining items: card, eraser, sponge. But maybe they value some of those. However, if their values are different, maybe they want some of the items I have. But since I don't know their values, I need to communicate.} \\
\\ 
\texttt{The rules say we can exchange up to 5 messages each. So I should start by sending a message to the other player to discuss possible allocations. Let me think about what to propose. Since I don't know their preferences, maybe I can suggest a division that allows them to take items I don't value. For example, if I take the two hooks and the flower, I can offer the other player the eraser and sponge. But the card is also there. Alternatively, maybe I can take one hook and the flower, and let them take the other hook and the eraser. But I need to make sure that the other player can get their desired items without overlapping with mine.} \\
\\ 
\texttt{Wait, the key is to make sure that the proposals are complementary. So if I propose taking certain items, the other player's proposal must not require any of those items. For example, if I take both hooks and the flower, the other player can't take any of those. But if they want the eraser or sponge, that's okay. But if they also want a hook or the flower, that would conflict. So maybe I should propose taking one hook and the flower, and then offer the other hook and the eraser to them. That way, if they value the} \\
            }
        }
    }
    & & \\ \\

    \theutterance \stepcounter{utterance}  
    & & \multicolumn{4}{p{0.6\linewidth}}{
        \cellcolor[rgb]{0.9,0.9,0.9}{
            \makecell[{{p{\linewidth}}}]{
                \texttt{\tiny{[GM$\rangle$P2]}}
                \texttt{<think>} \\
\texttt{Okay, let's try to figure this out. I need to divide the items between me and the other player. The items are 2 hooks, 1 flower, 1 card, 1 eraser, 1 sponge. My values are hook:3, flower:3, eraser:1, others 0. The other player's values are secret, so I don't know what they want. The goal is to maximize the total score for both of us, so we need to find a way to split the items where our proposals don't conflict.} \\
\\ 
\texttt{First, I should consider what I value most. Hooks and flowers are the highest for me. So ideally, I want to get as many of those as possible. But there are only two hooks and one flower. If I take both hooks and the flower, that's 3+3+3=9 points for me. But then the other player would have to take the remaining items: card, eraser, sponge. But maybe they value some of those. However, if their values are different, maybe they want some of the items I have. But since I don't know their values, I need to communicate.} \\
\\ 
\texttt{The rules say we can exchange up to 5 messages each. So I should start by sending a message to the other player to discuss possible allocations. Let me think about what to propose. Since I don't know their preferences, maybe I can suggest a division that allows them to take items I don't value. For example, if I take the two hooks and the flower, I can offer the other player the eraser and sponge. But the card is also there. Alternatively, maybe I can take one hook and the flower, and let them take the other hook and the eraser. But I need to make sure that the other player can get their desired items without overlapping with mine.} \\
\\ 
\texttt{Wait, the key is to make sure that the proposals are complementary. So if I propose taking certain items, the other player's proposal must not require any of those items. For example, if I take both hooks and the flower, the other player can't take any of those. But if they want the eraser or sponge, that's okay. But if they also want a hook or the flower, that would conflict. So maybe I should propose taking one hook and the flower, and then offer the other hook and the eraser to them. That way, if they value the} \\
            }
        }
    }
    & \\ \\

    \theutterance \stepcounter{utterance}  
    & & & \multicolumn{4}{p{0.6\linewidth}}{
        \cellcolor[rgb]{0.9,0.9,0.9}{
            \makecell[{{p{\linewidth}}}]{
                \texttt{\tiny{[GM$\langle$P2]}}
                \texttt{<think>} \\
\texttt{Okay, let's try to figure this out. I need to divide the items between me and the other player. The items are 2 hooks, 1 flower, 1 card, 1 eraser, 1 sponge. My values are hook:3, flower:3, eraser:1, others 0. The other player's values are secret, so I don't know what they want. The goal is to maximize the total score for both of us, so we need to find a way to split the items where our proposals don't conflict.} \\
\\ 
\texttt{First, I should consider what I value most. Hooks and flowers are the highest for me. So ideally, I want to get as many of those as possible. But there are only two hooks and one flower. If I take both hooks and the flower, that's 3+3+3=9 points for me. But then the other player would have to take the remaining items: card, eraser, sponge. But maybe they value some of those. However, if their values are different, maybe they want some of the items I have. But since I don't know their values, I need to communicate.} \\
\\ 
\texttt{The rules say we can exchange up to 5 messages each. So I should start by sending a message to the other player to discuss possible allocations. Let me think about what to propose. Since I don't know their preferences, maybe I can suggest a division that allows them to take items I don't value. For example, if I take the two hooks and the flower, I can offer the other player the eraser and sponge. But the card is also there. Alternatively, maybe I can take one hook and the flower, and let them take the other hook and the eraser. But I need to make sure that the other player can get their desired items without overlapping with mine.} \\
\\ 
\texttt{Wait, the key is to make sure that the proposals are complementary. So if I propose taking certain items, the other player's proposal must not require any of those items. For example, if I take both hooks and the flower, the other player can't take any of those. But if they want the eraser or sponge, that's okay. But if they also want a hook or the flower, that would conflict. So maybe I should propose taking one hook and the flower, and then offer the other hook and the eraser to them. That way, if they value the} \\
            }
        }
    }
     \\ \\

    \theutterance \stepcounter{utterance}  
    & & & \multicolumn{2}{p{0.3\linewidth}}{
        \cellcolor[rgb]{0.9,0.9,0.9}{
            \makecell[{{p{\linewidth}}}]{
                \texttt{\tiny{[GM$|$GM]}}
                \texttt{continue} \\
            }
        }
    }
    & & \\ \\

    \theutterance \stepcounter{utterance}  
    & & & \multicolumn{2}{p{0.3\linewidth}}{
        \cellcolor[rgb]{0.9,0.9,0.9}{
            \makecell[{{p{\linewidth}}}]{
                \texttt{\tiny{[GM$|$GM]}}
                \texttt{<think>} \\
\texttt{Okay, let's try to figure this out. I need to divide the items between me and the other player. The items are 2 hooks, 1 flower, 1 card, 1 eraser, 1 sponge. My values are hook:3, flower:3, eraser:1, others 0. The other player's values are secret, so I don't know what they want. The goal is to maximize the total score for both of us, so we need to find a way to split the items where our proposals don't conflict.} \\
\\ 
\texttt{First, I should consider what I value most. Hooks and flowers are the highest for me. So ideally, I want to get as many of those as possible. But there are only two hooks and one flower. If I take both hooks and the flower, that's 3+3+3=9 points for me. But then the other player would have to take the remaining items: card, eraser, sponge. But maybe they value some of those. However, if their values are different, maybe they want some of the items I have. But since I don't know their values, I need to communicate.} \\
\\ 
\texttt{The rules say we can exchange up to 5 messages each. So I should start by sending a message to the other player to discuss possible allocations. Let me think about what to propose. Since I don't know their preferences, maybe I can suggest a division that allows them to take items I don't value. For example, if I take the two hooks and the flower, I can offer the other player the eraser and sponge. But the card is also there. Alternatively, maybe I can take one hook and the flower, and let them take the other hook and the eraser. But I need to make sure that the other player can get their desired items without overlapping with mine.} \\
\\ 
\texttt{Wait, the key is to make sure that the proposals are complementary. So if I propose taking certain items, the other player's proposal must not require any of those items. For example, if I take both hooks and the flower, the other player can't take any of those. But if they want the eraser or sponge, that's okay. But if they also want a hook or the flower, that would conflict. So maybe I should propose taking one hook and the flower, and then offer the other hook and the eraser to them. That way, if they value the} \\
            }
        }
    }
    & & \\ \\

    \theutterance \stepcounter{utterance}  
    & & \multicolumn{4}{p{0.6\linewidth}}{
        \cellcolor[rgb]{0.9,0.9,0.9}{
            \makecell[{{p{\linewidth}}}]{
                \texttt{\tiny{[P1$\langle$GM]}}
                \texttt{<think>} \\
\texttt{Okay, let's try to figure this out. I need to divide the items between me and the other player. The items are 2 hooks, 1 flower, 1 card, 1 eraser, 1 sponge. My values are hook:3, flower:3, eraser:1, others 0. The other player's values are secret, so I don't know what they want. The goal is to maximize the total score for both of us, so we need to find a way to split the items where our proposals don't conflict.} \\
\\ 
\texttt{First, I should consider what I value most. Hooks and flowers are the highest for me. So ideally, I want to get as many of those as possible. But there are only two hooks and one flower. If I take both hooks and the flower, that's 3+3+3=9 points for me. But then the other player would have to take the remaining items: card, eraser, sponge. But maybe they value some of those. However, if their values are different, maybe they want some of the items I have. But since I don't know their values, I need to communicate.} \\
\\ 
\texttt{The rules say we can exchange up to 5 messages each. So I should start by sending a message to the other player to discuss possible allocations. Let me think about what to propose. Since I don't know their preferences, maybe I can suggest a division that allows them to take items I don't value. For example, if I take the two hooks and the flower, I can offer the other player the eraser and sponge. But the card is also there. Alternatively, maybe I can take one hook and the flower, and let them take the other hook and the eraser. But I need to make sure that the other player can get their desired items without overlapping with mine.} \\
\\ 
\texttt{Wait, the key is to make sure that the proposals are complementary. So if I propose taking certain items, the other player's proposal must not require any of those items. For example, if I take both hooks and the flower, the other player can't take any of those. But if they want the eraser or sponge, that's okay. But if they also want a hook or the flower, that would conflict. So maybe I should propose taking one hook and the flower, and then offer the other hook and the eraser to them. That way, if they value the} \\
\\ 
\\ 
\texttt{The time is up. It is now your turn to submit a secret proposal.} \\
            }
        }
    }
    & \\ \\

    \theutterance \stepcounter{utterance}  
    & \multicolumn{4}{p{0.6\linewidth}}{
        \cellcolor[rgb]{0.9,0.9,0.9}{
            \makecell[{{p{\linewidth}}}]{
                \texttt{\tiny{[P1$\rangle$GM]}}
                \texttt{<think>} \\
\texttt{Okay, let's try to figure this out. I need to divide the items between me and the other player. The items are 2 hooks, 1 flower, 1 card, 1 eraser, 1 sponge. My values are hook:3, flower:3, eraser:1, others 0. The other player's values are secret, so I don't know what they want. The goal is to maximize the total score for both of us, so we need to find a way to split the items where our proposals don't conflict.} \\
\\ 
\texttt{First, I should consider what I value most. Hooks and flowers are the highest for me. So ideally, I want to get as many of those as possible. But there are only two hooks and one flower. If I take both hooks and the flower, that's 3+3+3=9 points for me. But then the other player would have to take the remaining items: card, eraser, sponge. But maybe they value some of those. However, if their values are different, maybe they want some of the items I have. But since I don't know their values, I need to communicate.} \\
\\ 
\texttt{The rules say we can exchange up to 5 messages each. So I should start by sending a message to the other player to discuss possible allocations. Let me think about what to propose. Since I don't know their preferences, maybe I can suggest a division that allows them to take items I don't value. For example, if I take the two hooks and the flower, I can offer the other player the eraser and sponge. But the card is also there. Alternatively, maybe I can take one hook and the flower, and let them take the other hook and the eraser. But I need to make sure that the other player can get their desired items without overlapping with mine.} \\
\\ 
\texttt{Wait, the key is to make sure that the proposals are complementary. So if I propose taking certain items, the other player's proposal must not require any of those items. For example, if I take both hooks and the flower, the other player can't take any of those. But if they want the eraser or sponge, that's okay. But if they also want a hook or the flower, that would conflict. So maybe I should propose taking one hook and the flower, and then offer the other hook and the eraser to them. That way, if they value the} \\
            }
        }
    }
    & & \\ \\

    \theutterance \stepcounter{utterance}  
    & & & \multicolumn{2}{p{0.3\linewidth}}{
        \cellcolor[rgb]{0.9,0.9,0.9}{
            \makecell[{{p{\linewidth}}}]{
                \texttt{\tiny{[GM$|$GM]}}
                \texttt{continue} \\
            }
        }
    }
    & & \\ \\

    \theutterance \stepcounter{utterance}  
    & & & \multicolumn{2}{p{0.3\linewidth}}{
        \cellcolor[rgb]{0.9,0.9,0.9}{
            \makecell[{{p{\linewidth}}}]{
                \texttt{\tiny{[GM$|$GM]}}
                \texttt{player was instructed to make a proposal} \\
            }
        }
    }
    & & \\ \\

\end{supertabular}
}

\end{document}
